\documentclass[../../main.tex]{subfiles}
\begin{document}
\section{Whitney 嵌入定理}\label{section:ch06-Whitney嵌入定理}

我们应用 Sard 定理的第一个结果是证明每个光滑流形都可嵌入某 Euclid 空间。事实上,我们将证明每个带边或不带边光滑 $n$ 流形都微分同胚于 $\mathbb{R}^{2n+1}$ 的一个真嵌入子流形(带边或不带边)。

第一步是证明:若 $n$ 流形到 $\mathbb{R}^N$ 的单射浸入当 $N>2n+1$ 时,可改为到维数更低的 Euclid 空间的单射浸入。

\begin{lemma}\label{lemma:ch06-reduce-codimension}
设 $M\subseteq\mathbb{R}^N$ 是带边或不带边光滑 $n$ 维子流形。对任一 $v\in\mathbb{R}^N\setminus\mathbb{R}^{N-1}$,令 $\pi_v:\mathbb{R}^N\to\mathbb{R}^{N-1}$ 为以 $\mathbb{R}v$ 为核的投影(其中我们把 $\mathbb{R}^{N-1}$ 等同于 $\mathbb{R}^N$ 中最后一个坐标为零的点构成的子空间)。若 $N>2n+1$,则存在稠密集的向量 $v\in\mathbb{R}^N\setminus\mathbb{R}^{N-1}$,使 $\pi_v|_M$ 是 $M$ 到 $\mathbb{R}^{N-1}$ 的单射浸入。
\end{lemma}

\begin{proof}
为使 $\pi_v|_M$ 单射,当 $p,q$ 为 $M$ 中不同两点时, $p-q$ 永不平行于 $v$,这是必要且充分的。类似地,为使 $\pi_v|_M$ 是光滑浸入,对任一 $p\in M$, $T_pM$ 不含任何与 $\ker(\mathrm{d}\pi_v)_p$ 中非零向量,这是必要且充分的。由于 $\pi_v$ 是线性的,其微分就是同一线性映射(在通常的 $T_p\mathbb{R}^N\cong\mathbb{R}^N$ 等同下),故这一条件等价于要求 $T_pM$ 不含任何与 $v$ 平行的非零向量。

令 $\Delta_M\subseteq M\times M$ 表示闭集 $\Delta_M=\{(p,p):p\in M\}$(称为 $M\times M$ 的对角线),并令 $M_0\subseteq TM$ 表示闭集 $M_0=\{(p,0)\in TM:p\in M\}$( $M$ 各点处零向量之集)。考虑下面两个到实射影空间 $\mathbb{RP}^{N-1}$ 的映射:
\begin{align*}
\widetilde{\sigma}:(M\times M)\setminus\Delta_M&\to\mathbb{RP}^{N-1},&
\widetilde{\sigma}(p,q)&=[p-q],\\
\widetilde{\tau}:TM\setminus M_0&\to\mathbb{RP}^{N-1},&
\widetilde{\tau}(p,w)&=[w],
\end{align*}
其中方括号表示 $\mathbb{R}^N\setminus\{0\}$ 中向量作为 $\mathbb{RP}^{N-1}$ 中点的等价类。二者都光滑,因为它们是光滑映射与投影 $\mathbb{R}^N\setminus\{0\}\to\mathbb{RP}^{N-1}$ 的复合;而 $\pi_v|_M$ 为单射光滑浸入这一条件,恰为 $[v]$ 不在 $\widetilde{\sigma}$ 或 $\widetilde{\tau}$ 的像中这一条件。由于 $\widetilde{\sigma}$ 与 $\widetilde{\tau}$ 的定义域维数均为 $2n<N-1=\dim\mathbb{RP}^{N-1}$, Sard 定理的\cref{corollary:ch06-image-zero-dim} 表明每个映射的像零测,故其并也零测。于是等价类不在任一像中的向量之集稠密。
\end{proof}

通过反复应用上一引理,可得出结论:若 $n$ 流形容许到某 Euclid 空间的单射浸入,则它容许到 $\mathbb{R}^{2n+1}$ 的单射浸入。当 $M$ 紧时,由\cref{proposition:ch04-injective-immersion-embed}(c) 知此映射实际上是嵌入;但若 $M$ 不紧,我们需多做一点工作以确保单射浸入也是嵌入。

\begin{lemma}\label{lemma:ch06-proper-embedding-2n1}
设 $M$ 是带边或不带边光滑 $n$ 流形。若 $M$ 对某 $N$ 容许到 $\mathbb{R}^N$ 的光滑嵌入,则它容许到 $\mathbb{R}^{2n+1}$ 的真光滑嵌入。
\end{lemma}

\begin{proof}
在本证明中,给定一维线性子空间 $S\subseteq\mathbb{R}^N$ 与正数 $R$,我们把以 $S$ 为轴、半径为 $R$ 的管定义为 $\mathbb{R}^N$ 中开子集 $T_R(S)$,由到 $S$ 的距离小于 $R$ 的点组成:
\begin{align*}
T_R(S)=\{x\in\mathbb{R}^N:|x-y|<R\text{ 对某 }y\in S\}.
\end{align*}

设 $F:M\to\mathbb{R}^N$ 为任一光滑嵌入。令 $G:\mathbb{R}^N\to B^N$ 为微分同胚,并令 $f:M\to\mathbb{R}$ 为光滑穷竭函数(见\cref{proposition:ch02-exhaustion})。定义 $\widetilde{F}:M\to\mathbb{R}^N\times\mathbb{R}$ 为 $\widetilde{F}(p)=\bigl(G\circ F(p),f(p)\bigr)$。由于 $G\circ F$ 是嵌入,故 $\widetilde{F}$ 单射且 $\mathrm{d}\widetilde{F}_p$ 对每个 $p$ 单射,从而 $\widetilde{F}$ 是单射浸入。它固有,因为任一紧集的原像是某 $f^{-1}((-\infty,c])$ 的闭子集,而后者紧;故由\cref{proposition:ch04-injective-immersion-embed}(b) 知 $\widetilde{F}$ 是光滑嵌入。由构造, $\widetilde{F}$ 的像含于管 $B^N\times\mathbb{R}$ 中。

此后(把 $N+1$ 换成 $N$),我们假设 $M$ 容许到 $\mathbb{R}^N$ 的真光滑嵌入,其值落在某管 $T_R(S)$ 中(\cref{figure:ch06-reduce-codimension})。把 $M$ 与其像等同,可把 $M$ 视为含于该管中的 $\mathbb{R}^N$ 的真嵌入子流形。

若 $N>2n+1$, \cref{lemma:ch06-reduce-codimension} 表明存在 $v\in\mathbb{R}^N\setminus\mathbb{R}^{N-1}$,使 $\pi_v|_M$ 是 $M$ 到 $\mathbb{R}^{N-1}$ 的单射浸入。此外,我们可选取 $v$ 使它不在子空间 $S$ 中;于是 $\pi_v(S)$ 是 $\mathbb{R}^{N-1}$ 的一维子空间,且 $\pi_v(M)$ 含于 $\pi_v(S)$ 周围的某管中,因为 $\pi_v$ 是有界线性映射。我们将证明 $\pi_v|_M$ 固有,故由\cref{proposition:ch04-injective-immersion-embed}(b) 知它是嵌入。

设 $K\subseteq\mathbb{R}^{N-1}$ 为紧集。则 $K$ 含于以 $0$ 为中心、某半径 $R_1$ 的开球中。对任一 $x\in\pi_v^{-1}(K)$,存在某 $c\in\mathbb{R}$ 使 $\pi_v(x)=x-cv$;由于 $|\pi_v(x)|<R_1$,这意味着 $x$ 在以 $v$ 张成的直线 $\mathbb{R}v$ 为轴、半径 $R_1$ 的管中。于是 $M\cap\pi_v^{-1}(K)$ 含于两个管中,一个绕 $S$,另一个绕 $\mathbb{R}v$。一个简单的几何论证表明:当两管的轴不平行时,其交有界,故 $M\cap\pi_v^{-1}(K)$ 紧。于是 $\pi_v|_M$ 固有,这蕴涵 $\pi_v(M)$ 是 $\mathbb{R}^{N-1}$ 中含于某管的真嵌入子流形。现在可迭代此论证,直到得到 $M$ 到 $\mathbb{R}^{2n+1}$ 的真光滑嵌入。
\end{proof}

\begin{figure}[ht]
\centering
\includegraphics[width=0.8\textwidth]{fig-06-5.png}
\caption{降低逆紧嵌入的余维数}\label{figure:ch06-reduce-codimension}
\end{figure}

\begin{theorem}[Whitney Embedding Theorem]\label{theorem:ch06-whitney-embedding}
每个带边或不带边光滑 $n$ 流形都容许到 $\mathbb{R}^{2n+1}$ 的真光滑嵌入。
\end{theorem}

\begin{proof}
设 $M$ 是带边或不带边光滑 $n$ 流形。由\cref{lemma:ch06-proper-embedding-2n1},只需证明 $M$ 容许到某 Euclid 空间的光滑嵌入。

先设 $M$ 紧。在此情形, $M$ 容许有限覆盖 $\{B_1,\dots,B_m\}$,其中每个 $B_i$ 是正则坐标球或正则坐标半球。这意味着对每个 $i$,存在坐标域 $B_i'\supseteq\overline{B_i}$ 与光滑坐标映射 $\varphi_i:B_i'\to\mathbb{R}^n$,它限制为从 $\overline{B_i}$ 到 $\mathbb{R}^n$ 的紧子集上的微分同胚。对每个 $i$,令 $\rho_i:M\to\mathbb{R}$ 为光滑鼓包函数,在 $\overline{B_i}$ 上等于 $1$,支集含于 $B_i'$。定义光滑映射 $F:M\to\mathbb{R}^{nm+m}$ 为
\begin{align*}
F(p)=\bigl(\rho_1(p)\varphi_1(p),\dots,\rho_m(p)\varphi_m(p),\rho_1(p),\dots,\rho_m(p)\bigr),
\end{align*}
其中如通常一样, $\rho_i\varphi_i$ 在 $\rho_i$ 支集之外延拓为零。我们将证明 $F$ 是单射光滑浸入;由于 $M$ 紧,这蕴涵 $F$ 是光滑嵌入。

为看 $F$ 单射,设 $F(p)=F(q)$。由于集族 $\{B_i\}$ 覆盖 $M$,存在某 $i$ 使 $p\in B_i$。于是 $\rho_i(p)=1$,而 $\rho_i(q)=\rho_i(p)=1$ 这一事实蕴涵 $q\in\operatorname{supp}\rho_i\subseteq B_i'$,且
\begin{align*}
\varphi_i(q)=\rho_i(q)\varphi_i(q)=\rho_i(p)\varphi_i(p)=\varphi_i(p).
\end{align*}
由于 $\varphi_i$ 在 $B_i'$ 上单射,故 $p=q$。

其次,为看 $F$ 是光滑浸入,任取 $p\in M$,并选 $i$ 使 $p\in B_i$。由于 $\rho_i$ 在 $p$ 的某邻域上等于 $1$,有 $\mathrm{d}(\rho_i\varphi_i)_p=\mathrm{d}(\varphi_i)_p$,它单射。由此易知 $\mathrm{d}F_p$ 单射。于是 $F$ 是单射光滑浸入,因而是嵌入。

现设 $M$ 非紧。令 $f:M\to\mathbb{R}$ 为光滑穷竭函数。Sard 定理表明:对每个非负整数 $i$,存在 $f$ 的正则值 $a_i,b_i$,使 $i<a_i<b_i<i+1$。定义子集 $D_i,E_i\subseteq M$ 为
\begin{align*}
D_0=f^{-1}((-\infty,1]),\qquad E_0=f^{-1}((-\infty,a_1]),
\end{align*}
\begin{align*}
D_i=f^{-1}([i,i+1]),\qquad E_i=f^{-1}([b_{i-1},a_{i+1}]),\quad i\geqslant 1.
\end{align*}
由\cref{proposition:ch05-proper-closed},每个 $E_i$ 是紧正则区域。我们有 $D_i\subseteq\operatorname{Int}E_i$, $M=\bigcup_i D_i$,且 $E_i\cap E_j=\varnothing$ 除非 $j=i-1,i,i+1$。证明的第一部分表明每个 $E_i$ 都有到某 Euclid 空间的光滑嵌入,故由\cref{lemma:ch06-proper-embedding-2n1} 存在嵌入 $\varphi_i:E_i\to\mathbb{R}^{2n+1}$。对每个 $i$,令 $\rho_i:M\to\mathbb{R}$ 为光滑鼓包函数,在 $D_i$ 的某邻域上等于 $1$,支集含于 $\operatorname{Int}E_i$,并定义 $F:M\to\mathbb{R}^{2n+1}\times\mathbb{R}^{2n+1}\times\mathbb{R}$ 为
\begin{align*}
F(p)=\left(\sum_{i\text{ 偶}}\rho_i(p)\varphi_i(p),\ \sum_{i\text{ 奇}}\rho_i(p)\varphi_i(p),\ f(p)\right).
\end{align*}
则 $F$ 光滑,因为在每点附近每个和中只有一项非零;且 $F$ 固有,因为 $f$ 固有。我们将证明 $F$ 也是单射光滑浸入,因而是光滑嵌入。

设 $F(p)=F(q)$。则 $p\in D_j$ 对某 $j$,而 $f(q)=f(p)$ 蕴涵 $q\in D_j$。完全仿上紧情形的论证,得出 $p=q$。

现任取 $p\in M$,并选 $j$ 使 $p\in D_j$。则 $\rho_j$ 在 $p$ 的某邻域上等于 $1$。设 $j$ 为奇,则对该邻域中所有 $q$,有
\begin{align*}
F(q)=\bigl(\varphi_j(q),\cdot,\cdot\bigr),
\end{align*}
这蕴涵 $\mathrm{d}F_p$ 单射。当 $j$ 为偶时类似。
\end{proof}

\begin{corollary}\label{corollary:ch06-diffeomorphic-embedded}
每个带边或不带边光滑 $n$ 流形都微分同胚于 $\mathbb{R}^{2n+1}$ 的一个真嵌入子流形(带边或不带边)。
\end{corollary}

\begin{corollary}\label{corollary:ch06-approximate-embedding}
设 $M$ 是带边或不带边紧光滑 $n$ 流形。若 $N\geqslant 2n+1$,则每个从 $M$ 到 $\mathbb{R}^N$ 的光滑映射都可被嵌入一致逼近。
\end{corollary}

\begin{proof}
设 $N\geqslant 2n+1$,并令 $f:M\to\mathbb{R}^N$ 为光滑映射。由 Whitney 嵌入定理,存在光滑嵌入 $F:M\to\mathbb{R}^{2n+1}$。映射 $G=f\times F:M\to\mathbb{R}^N\times\mathbb{R}^{2n+1}$ 也是光滑嵌入,而 $f$ 等于复合 $\pi\circ G$,其中 $\pi:\mathbb{R}^N\times\mathbb{R}^{2n+1}\to\mathbb{R}^N$ 是投影。令 $\widetilde{M}=G(M)\subseteq\mathbb{R}^N\times\mathbb{R}^{2n+1}$。\cref{lemma:ch06-reduce-codimension} 表明存在任意接近 $e_{N+2n+1}=(0,\dots,0,1)$ 的向量 $v_{N+2n+1}\in\mathbb{R}^N\times\mathbb{R}^{2n+1}$,使 $\pi_{v_{N+2n+1}}|_{\widetilde{M}}$ 是嵌入。这蕴涵 $\pi_{v_{N+2n+1}}$ 任意接近 $\pi_{e_{N+2n+1}}$。迭代此过程,得到向量 $v_{N+2n+1},v_{N+2n},\dots,v_{N+1}$,使 $\pi_{v_{N+1}}\circ\cdots\circ\pi_{v_{N+2n+1}}$ 限制为 $\widetilde{M}$ 的一个嵌入,它任意接近 $\pi_{e_{N+1}}\circ\cdots\circ\pi_{e_{N+2n+1}}=\pi$,从而 $\pi_{v_{N+1}}\circ\cdots\circ\pi_{v_{N+2n+1}}\circ G$ 是 $M$ 到 $\mathbb{R}^N$ 的嵌入,它任意接近 $f$。
\end{proof}

若我们只要求浸入而非嵌入,则维数可降低一。

\begin{theorem}[Whitney Immersion Theorem]\label{theorem:ch06-whitney-immersion}
每个带边或不带边光滑 $n$ 流形都容许到 $\mathbb{R}^{2n}$ 的光滑浸入。
\end{theorem}

\begin{proof}
$\partial M=\varnothing$ 的情形见习题 6-2,一般情形见习题 9-14。
\end{proof}

\cref{theorem:ch06-whitney-embedding} 由 Hassler Whitney 于 1936 年(在边界为空的情形)首先证明 [Whi36],回答了一个自抽象流形概念首次提出以来一直困扰数学家的问题:是否存在不同胚于 Euclid 空间嵌入子流形的抽象光滑流形?现在我们知道不存在。

尽管我们证明的这一版本嵌入定理对我们的目的已足够,但值得注意的是,八年后,利用代数拓扑中更为精细的技术,Whitney 得到了如下改进 [Whi44a, Whi44b]。

\begin{theorem}[Strong Whitney Embedding Theorem]\label{theorem:ch06-strong-embedding}
若 $n>0$,每个光滑 $n$ 流形都容许到 $\mathbb{R}^{2n}$ 的光滑嵌入。
\end{theorem}

\begin{theorem}[Strong Whitney Immersion Theorem]\label{theorem:ch06-strong-immersion}
若 $n>1$,每个光滑 $n$ 流形都容许到 $\mathbb{R}^{2n-1}$ 的光滑浸入。
\end{theorem}

由于这些结果, \cref{theorem:ch06-whitney-embedding} 与\cref{theorem:ch06-whitney-immersion} 有时称为弱 Whitney 嵌入与浸入定理。

事实上,即使是强 Whitney 定理在所有维数下也未必最优。1985 年,Ralph Cohen 证明每个紧光滑 $n$ 流形都可浸入 $\mathbb{R}^{2n-\alpha(n)}$,其中 $\alpha(n)$ 是 $n$ 的二进制表达式中 $1$ 的个数。例如,每个 $3$ 流形可浸入 $\mathbb{R}^4$,而 $4$ 流形需要 $\mathbb{R}^7$。这个结果在每一维数下都是最优的。另一方面,最优嵌入维数只在某些维数下已知。例如, Whitney 的维数 $2n$ 对 $n=1$ 与 $n=2$ 的流形最优,但 C.T.C. Wall 于 1965 年证明 [Wal65],每个 $3$ 流形都可嵌入 $\mathbb{R}^5$。关于最新进展及其参考文献的一个好综述见 [Osb82]。

\end{document}